% Physical pages: 81--97 (printed pages 58--74).
% Figures: none.
% Uncertainties: none.

\subsection{The Poincar\'e Algebra}

As we saw in Section 2.2, much of the information about any Lie symmetry group
is contained in properties of the group elements near the identity. For the
Poincar\'e group, the identity is the transformation
$\Lambda^\mu{}_{\nu}=\delta^\mu{}_{\nu}$, $a^\mu=0$, so we want to study those
transformations with
\begin{equation}
\Lambda^\mu{}_{\nu}=\delta^\mu{}_{\nu}+\omega^\mu{}_{\nu}\,,
\qquad
a^\mu=\epsilon^\mu\,,
\tag{2.4.1}\label{eq:2.4.1}
\end{equation}
both $\omega^\mu{}_{\nu}$ and $\epsilon^\mu$ being taken infinitesimal. The
Lorentz condition \eqref{eq:2.3.5}

reads here
\[
\begin{aligned}
\eta_{\rho\sigma}
  &= \eta_{\mu\nu}
     (\delta^\mu_{\rho}+\omega^\mu{}_{\rho})
     (\delta^\nu_{\sigma}+\omega^\nu{}_{\sigma}) \\
  &= \eta_{\rho\sigma}+\omega_{\sigma\rho}
     +\omega_{\rho\sigma}+O(\omega^2)\,.
\end{aligned}
\]
We are here using the convention, to be used throughout this book, that
indices may be lowered or raised by contraction with $\eta_{\mu\nu}$ or
$\eta^{\mu\nu}$
\[
\omega_{\sigma\rho}=\eta_{\mu\sigma}\omega^\mu{}_{\rho},
\qquad
\omega^\mu{}_{\rho}=\eta^{\mu\sigma}\omega_{\sigma\rho}\,.
\]
Keeping only the terms of first order in $\omega$ in the Lorentz condition
\eqref{eq:2.3.5}, we see that this condition now reduces to the antisymmetry of
$\omega_{\mu\nu}$
\begin{equation}
\omega_{\mu\nu}=-\omega_{\nu\mu}\,.
\tag{2.4.2}\label{eq:2.4.2}
\end{equation}
An antisymmetric second-rank tensor in four dimensions has $(4\cdot3)/2=6$
independent components, so including the four components of $\epsilon^\mu$,
an inhomogeneous Lorentz transformation is described by $6+4=10$
parameters.

Since $U(\mathbf{1},0)$ carries any ray into itself, it must be proportional to the unit
operator,\footnote{In the absence of superselection rules, the possibility that
the proportionality factor may depend on the state on which $U(\mathbf{1},0)$ acts can
be ruled out by the same reasoning that we used in Section 2.2 to rule out the
possibility that the phases in projective representations of symmetry groups
may depend on the states on which the symmetries act. Where superselection
rules apply, it may be necessary to redefine $U(\mathbf{1},0)$ by phase factors that
depend on the sector on which it acts.} and by a choice of phase may be made
equal to it. For an infinitesimal Lorentz transformation \eqref{eq:2.4.1},
$U(\mathbf{1}+\omega,\epsilon)$ must then equal $\mathbf{1}$ plus terms linear in $\omega_{\rho\sigma}$
and $\epsilon_\rho$. We write this as
\begin{equation}
U(\mathbf{1}+\omega,\epsilon)
=\mathbf{1}+\frac{i}{2}\omega_{\rho\sigma}J^{\rho\sigma}
-i\epsilon_\rho P^\rho+\cdots\,.
\tag{2.4.3}\label{eq:2.4.3}
\end{equation}
Here $J^{\rho\sigma}$ and $P^\rho$ are $\omega$- and $\epsilon$-independent
operators, and the dots denote terms of higher order in $\omega$ and/or
$\epsilon$. In order for $U(\mathbf{1}+\omega,\epsilon)$ to be unitary, the operators
$J^{\rho\sigma}$ and $P^\rho$ must be Hermitian
\begin{equation}
J^{\rho\sigma\dagger}=J^{\rho\sigma}\,,
\qquad P^{\rho\dagger}=P^\rho\,.
\tag{2.4.4}\label{eq:2.4.4}
\end{equation}
Since $\omega_{\rho\sigma}$ is antisymmetric, we can take its coefficient
$J^{\rho\sigma}$ to be antisymmetric also
\begin{equation}
J^{\rho\sigma}=-J^{\sigma\rho}\,.
\tag{2.4.5}\label{eq:2.4.5}
\end{equation}
As we shall see, $P^1$, $P^2$, and $P^3$ are the components of the momentum
operators, $J^{23}$, $J^{31}$, and $J^{12}$ are the components of the angular
momentum vector, and $P^0$ is the energy operator, or
Hamiltonian.\footnote{We will see that this identification of the
angular-momentum generators is forced on us by the commutation relations of
the $J^{\mu\nu}$. On the other hand, the commutation relations do not allow us
to distinguish between $P^\mu$ and $-P^\mu$, so the sign for the
$\epsilon_\rho P^\rho$ term in \eqref{eq:2.4.3} is a matter of convention. The
consistency of the choice in \eqref{eq:2.4.3} with the usual definition of the
Hamiltonian $P^0$ is shown in Section 3.1.}

We now examine the Lorentz transformation properties of $J^{\rho\sigma}$ and
$P^\rho$. We consider the product
\[
U(\Lambda,a)\,U(\mathbf{1}+\omega,\epsilon)\,U^{-1}(\Lambda,a)\,,
\]
where $\Lambda^\mu{}_{\nu}$ and $a^\mu$ are here the parameters of a new
transformation, unrelated to $\omega$ and $\epsilon$. According to Eq.
\eqref{eq:2.3.11}, the product $U(\Lambda^{-1},-\Lambda^{-1}a)U(\Lambda,a)$ equals
$U(\mathbf{1},0)$, so $U(\Lambda^{-1},-\Lambda^{-1}a)$ is the inverse of $U(\Lambda,a)$.
It follows then from \eqref{eq:2.3.11} that
\begin{equation}
\begin{split}
&U(\Lambda,a)U(\mathbf{1}+\omega,\epsilon)U^{-1}(\Lambda,a)\\
&\qquad={}
U\!\left(\Lambda(\mathbf{1}+\omega)\Lambda^{-1},
\Lambda\epsilon-\Lambda\omega\Lambda^{-1}a\right)\,.
\end{split}
\tag{2.4.6}\label{eq:2.4.6}
\end{equation}
To first order in $\omega$ and $\epsilon$, we have then
\begin{equation}
\begin{split}
&U(\Lambda,a)
\left[\frac12\omega_{\rho\sigma}J^{\rho\sigma}
-\epsilon_\rho P^\rho\right]U^{-1}(\Lambda,a)\\
&\qquad={}
\frac12(\Lambda\omega\Lambda^{-1})_{\mu\nu}J^{\mu\nu}
-(\Lambda\epsilon-\Lambda\omega\Lambda^{-1}a)_\mu P^\mu\,.
\end{split}
\tag{2.4.7}\label{eq:2.4.7}
\end{equation}
Equating coefficients of $\omega_{\rho\sigma}$ and $\epsilon_\rho$ on both
sides of this equation (and using \eqref{eq:2.3.10}), we find
\begin{equation}
\begin{aligned}
U(\Lambda,a)J^{\rho\sigma}U^{-1}(\Lambda,a)
&=\Lambda_\mu{}^\rho\Lambda_\nu{}^\sigma\left(J^{\mu\nu}-a^\mu P^\nu+a^\nu P^\mu\right),
\end{aligned}
\tag{2.4.8}\label{eq:2.4.8}
\end{equation}
\begin{equation}
U(\Lambda,a)P^\rho U^{-1}(\Lambda,a)
=\Lambda_\mu{}^\rho P^\mu\,.
\tag{2.4.9}\label{eq:2.4.9}
\end{equation}
For homogeneous Lorentz transformations (with $a^\mu=0$), these
transformation rules simply say that $J^{\mu\nu}$ is a tensor and $P^\mu$ is a
vector. For pure translations (with $\Lambda^\mu{}_{\nu}=\delta^\mu{}_{\nu}$),
they tell us that $P^\rho$ is translation-invariant, but $J^{\rho\sigma}$ is
not. In particular, the change of the space-space components of
$J^{\rho\sigma}$ under a spatial translation is just the usual change of the
angular momentum under a change of the origin relative to which the angular
momentum is calculated.

Next, let's apply rules \eqref{eq:2.4.8}, \eqref{eq:2.4.9} to a transformation that is itself
infinitesimal, i.e., $\Lambda^\mu{}_{\nu}=\delta^\mu{}_{\nu}+\omega'^\mu{}_{\nu}$
and $a^\mu=\epsilon'^\mu$, with infinitesimals $\omega'^\mu{}_{\nu}$ and
$\epsilon'^\mu$ unrelated to the previous $\omega$ and $\epsilon$. Using Eq.
\eqref{eq:2.4.3}, and keeping only terms of first order in $\omega'^\mu{}_{\nu}$ and
$\epsilon'^\mu$, Eqs.~\eqref{eq:2.4.8} and \eqref{eq:2.4.9} now become
\begin{equation}
i\left[\frac12\omega'_{\mu\nu}J^{\mu\nu}-\epsilon'_\mu P^\mu,
J^{\rho\sigma}\right]
=\omega'^\rho{}_{\mu}J^{\mu\sigma}
+\omega'^\sigma{}_{\nu}J^{\rho\nu}
-\epsilon'^\rho P^\sigma+\epsilon'^\sigma P^\rho\,,
\tag{2.4.10}\label{eq:2.4.10}
\end{equation}
\begin{equation}
i\left[\frac12\omega'_{\mu\nu}J^{\mu\nu}-\epsilon'_\mu P^\mu,
P^\rho\right]
=\omega'^\rho{}_{\mu}P^\mu\,.
\tag{2.4.11}\label{eq:2.4.11}
\end{equation}
Equating coefficients of $\omega'_{\mu\nu}$ and $\epsilon'_\mu$ on both sides
of these equations, we find the commutation rules
\begin{equation}
i[J^{\mu\nu},J^{\rho\sigma}]
=\eta^{\nu\rho}J^{\mu\sigma}-\eta^{\mu\rho}J^{\nu\sigma}
-\eta^{\nu\sigma}J^{\mu\rho}+\eta^{\mu\sigma}J^{\nu\rho}\,,
\tag{2.4.12}\label{eq:2.4.12}
\end{equation}
\begin{equation}
i[P^\mu,J^{\rho\sigma}]
=\eta^{\mu\rho}P^\sigma-\eta^{\mu\sigma}P^\rho\,,
\tag{2.4.13}\label{eq:2.4.13}
\end{equation}
\begin{equation}
[P^\mu,P^\rho]=0\,.
\tag{2.4.14}\label{eq:2.4.14}
\end{equation}
This is the Lie algebra of the Poincar\'e group.

In quantum mechanics a special role is played by those operators that are
\emph{conserved}, i.e., that commute with the energy operator $H=P^0$.
Inspection of Eqs.~\eqref{eq:2.4.13} and \eqref{eq:2.4.14} shows that these are the momentum
three-vector
\begin{equation}
\mathbf{P}=\{P^1,P^2,P^3\}
\tag{2.4.15}\label{eq:2.4.15}
\end{equation}
and the angular-momentum three-vector
\begin{equation}
\mathbf{J}=\{J^{23},J^{31},J^{12}\}
\tag{2.4.16}\label{eq:2.4.16}
\end{equation}
and, of course, the energy $P^0$ itself. The remaining generators form what is
called the `boost' three-vector
\begin{equation}
\mathbf{K}=\{J^{10},J^{20},J^{30}\}\,.
\tag{2.4.17}\label{eq:2.4.17}
\end{equation}
These are \emph{not} conserved, which is why we do not use the eigenvalues of
$\mathbf{K}$ to label physical states. In a three-dimensional notation, the
commutation relations \eqref{eq:2.4.12}, \eqref{eq:2.4.13}, \eqref{eq:2.4.14} may be written
\begin{align}
[J_i,J_j]&=i\epsilon_{ijk}J_k\,, \tag{2.4.18}\label{eq:2.4.18}\\
[J_i,K_j]&=i\epsilon_{ijk}K_k\,, \tag{2.4.19}\label{eq:2.4.19}\\
[K_i,K_j]&=-i\epsilon_{ijk}J_k\,, \tag{2.4.20}\label{eq:2.4.20}\\
[J_i,P_j]&=i\epsilon_{ijk}P_k\,, \tag{2.4.21}\label{eq:2.4.21}\\
[K_i,P_j]&=iH\delta_{ij}\,, \tag{2.4.22}\label{eq:2.4.22}\\
[J_i,H]&=[P_i,H]=[H,H]=0\,, \tag{2.4.23}\label{eq:2.4.23}\\
[K_i,H]&=iP_i\,, \tag{2.4.24}\label{eq:2.4.24}
\end{align}
where $i,j,k$, etc. run over the values 1, 2, and 3, and $\epsilon_{ijk}$ is
the totally antisymmetric quantity with $\epsilon_{123}=+1$. The commutation
relation \eqref{eq:2.4.18} will be recognized as that of the angular-momentum operator.

The pure translations $T(\mathbf{1},a)$ form a subgroup of the Poincar\'e
group with a group multiplication rule given by \eqref{eq:2.3.7} as
\begin{equation}
T(\mathbf{1},\bar a)T(\mathbf{1},a)=T(\mathbf{1},\bar a+a)\,.
\tag{2.4.25}\label{eq:2.4.25}
\end{equation}
This is additive in the same sense as \eqref{eq:2.2.24}, so by using \eqref{eq:2.4.3} and
repeating the same arguments that led to \eqref{eq:2.2.26}, we find that finite
translations are represented on the physical Hilbert space by
\begin{equation}
U(\mathbf{1},a)=\exp(-iP^\mu a_\mu)\,.
\tag{2.4.26}\label{eq:2.4.26}
\end{equation}
In exactly the same way, we can show that a rotation $R_{\boldsymbol\theta}$ by
an angle $|\boldsymbol\theta|$ around the direction of $\boldsymbol\theta$ is
represented on the physical Hilbert space by
\begin{equation}
U(R_{\boldsymbol\theta},0)=\exp(i\mathbf{J}\cdot\boldsymbol\theta)\,.
\tag{2.4.27}\label{eq:2.4.27}
\end{equation}

It is interesting to compare the Poincar\'e algebra with the Lie algebra of the
symmetry group of Newtonian mechanics, the Galilean group. We could derive this
algebra by starting with the transformation rules of the Galilean group and
then following the same procedure that was used here to derive the Poincar\'e
algebra. However, since we already have Eqs.~\eqref{eq:2.4.18}--\eqref{eq:2.4.24}, it is easier
to obtain the Galilean algebra as a low-velocity limit of the Poincar\'e algebra,
by what is known as an \emph{In\"on\"u--Wigner contraction}~[\hyperref[ref:2.4]{4},\hyperref[ref:2.5]{5}].
For a system of particles of typical mass $m$ and typical velocity $v$, the
momentum and the angular-momentum operators are expected to be of order
$J\sim1$, $P\sim mv$. On the other hand, the energy operator is $H=M+W$ with a
total mass $M$ and non-mass energy $W$ (kinetic plus potential) of order
$M\sim m$, $W\sim mv^2$. Inspection of Eqs.~\eqref{eq:2.4.18}--\eqref{eq:2.4.24} shows that these
commutation relations have a limit for $v\ll1$ of the form
\[
\begin{gathered}
[J_i,J_j]=i\epsilon_{ijk}J_k\,,\qquad
[J_i,K_j]=i\epsilon_{ijk}K_k\,,\qquad
[K_i,K_j]=0\,,\\
[J_i,P_j]=i\epsilon_{ijk}P_k\,,\qquad
[K_i,P_j]=iM\delta_{ij}\,,\\
[J_i,W]=[P_i,W]=0\,,\qquad [K_i,W]=iP_i\,,\\
[J_i,M]=[P_i,M]=[K_i,M]=[W,M]=0\,,
\end{gathered}
\]
with $K$ of order $1/v$. Note that the product of a translation
$\mathbf{x}\to\mathbf{x}+\mathbf{a}$ and a `boost'
$\mathbf{x}\to\mathbf{x}+\mathbf{v}t$ should be the transformation
$\mathbf{x}\to\mathbf{x}+\mathbf{v}t+\mathbf{a}$, but this is not true for the
action of these operators on Hilbert space:
\[
\exp(i\mathbf{K}\cdot\mathbf{v})\exp(-i\mathbf{P}\cdot\mathbf{a})
=\exp(iM\mathbf{a}\cdot\mathbf{v}/2)
 \exp\!\left(i(\mathbf{K}\cdot\mathbf{v}-\mathbf{P}\cdot\mathbf{a})\right).
\]
The appearance of the phase factor $\exp(iM\mathbf{a}\cdot\mathbf{v}/2)$ shows
that this is a projective representation, with a superselection rule forbidding
the superposition of states of different mass. In this respect, the mathematics
of the Poincar\'e group is simpler than that of the Galilean group. However,
there is nothing to prevent us from formally enlarging the Galilean group, by
adding one more generator to its Lie algebra, which commutes with all the other
generators, and whose eigenvalues are the masses of the various states. In this
case physical states provide an ordinary rather than a projective
representation of the expanded symmetry group. The difference appears to be a
mere matter of notation, except that with this reinterpretation of the Galilean
group there is no need for a mass superselection rule.

\subsection{One-Particle States}

We now consider the classification of one-particle states according to their
transformation under the Poincar\'e group.

The components of the energy-momentum four-vector all commute with each other,
so it is natural to express physical state-vectors in terms of eigenvectors of
the four-momentum. Introducing a label $\sigma$ to denote all other degrees of
freedom, we thus consider state-vectors $\ket{p,\sigma}$ with
\begin{equation}
P^\mu\ket{p,\sigma}=p^\mu\ket{p,\sigma}\,.
\tag{2.5.1}\label{eq:2.5.1}
\end{equation}
For general states, consisting for instance of several unbound particles, the
label $\sigma$ would have to be allowed to include continuous as well as
discrete labels. We take as part of the definition of a \emph{one-particle}
state, that the label $\sigma$ is purely discrete, and will limit ourselves
here to that case. (However, a specific bound state of two or more particles,
such as the lowest state of the hydrogen atom, is to be considered as a
one-particle state. It is not an \emph{elementary} particle, but the distinction
between composite and elementary particles is of no relevance here.)

Eqs.~\eqref{eq:2.5.1} and \eqref{eq:2.4.26} tell us how the states $\ket{p,\sigma}$ transform
under translations:
\[
U(\mathbf{1},a)\ket{p,\sigma}=e^{-ip\cdot a}\ket{p,\sigma}\,.
\]
We must now consider how these states transform under homogeneous Lorentz
transformations.

Using \eqref{eq:2.4.9}, we see that the effect of operating on $\ket{p,\sigma}$ with a
quantum homogeneous Lorentz transformation $U(\Lambda,0)\equiv U(\Lambda)$ is
to produce an eigenvector of the four-momentum with eigenvalue $\Lambda p$
\begin{equation}
\begin{aligned}
P^\mu U(\Lambda)\ket{p,\sigma}
&=U(\Lambda)\left[U^{-1}(\Lambda)P^\mu U(\Lambda)\right]\ket{p,\sigma}\\
&=U(\Lambda)(\Lambda^{-1\mu}{}_{\rho}P^\rho)\ket{p,\sigma}\\
&=\Lambda^\mu{}_{\rho}p^\rho U(\Lambda)\ket{p,\sigma}\,.
\end{aligned}
\tag{2.5.2}\label{eq:2.5.2}
\end{equation}
Hence $U(\Lambda)\ket{p,\sigma}$ must be a linear combination of the
state-vectors $\ket{\Lambda p,\sigma'}$:
\begin{equation}
U(\Lambda)\ket{p,\sigma}
=\sum_{\sigma'}C_{\sigma'\sigma}(\Lambda,p)
\ket{\Lambda p,\sigma'}\,.
\tag{2.5.3}\label{eq:2.5.3}
\end{equation}
In general, it may be possible by using suitable linear combinations of the
$\ket{p,\sigma}$ to choose the $\sigma$ labels in such a way that the matrix
$C_{\sigma'\sigma}(\Lambda,p)$ is block-diagonal; in other words, so that the
$\ket{p,\sigma}$ with $\sigma$ within any one block by themselves furnish a
representation of the Poincar\'e group. It is natural to identify
the states of a specific particle type with the components of a representation
of the Poincar\'e group which is irreducible, in the sense that it
cannot be further decomposed in this way.\footnote{Of course, different
particle species may correspond to representations that are isomorphic, i.e.,
that have matrices $C_{\sigma'\sigma}(\Lambda,p)$ that are either identical, or
identical up to a similarity transformation. In some cases it may be convenient
to define particle types as irreducible representations of larger groups that
contain the inhomogeneous proper orthochronous Lorentz group as a subgroup; for
instance, as we shall see, for massless particles whose interactions respect
the symmetry of space inversion it is customary to treat all the components of
an irreducible representation of the Poincar\'e group including
space inversion as a single particle type.} Our task now is to work out the
structure of the coefficients $C_{\sigma'\sigma}(\Lambda,p)$ in irreducible
representations of the Poincar\'e group.

For this purpose, note that the only functions of $p^\mu$ that are left
invariant by all proper orthochronous Lorentz transformations
$\Lambda^\mu{}_{\nu}$, are the invariant square
$p^2\equiv\eta_{\mu\nu}p^\mu p^\nu$, and for $p^2\leq0$, also the sign of
$p^0$. Hence, for each value of $p^2$, and (for $p^2\leq0$) each sign of
$p^0$, we can choose a `standard' four-momentum, say $k^\mu$, and express any
$p^\mu$ of this class as
\begin{equation}
p^\mu=L^\mu{}_{\nu}(p)k^\nu\,,
\tag{2.5.4}\label{eq:2.5.4}
\end{equation}
where $L^\mu{}_{\nu}$ is some standard Lorentz transformation that depends on
$p^\mu$, and also implicitly on our choice of the standard $k^\mu$. We can
then \emph{define} the states $\ket{p,\sigma}$ of momentum $p$ by
\begin{equation}
\ket{p,\sigma}\equiv N(p)U(L(p))\ket{k,\sigma}\,,
\tag{2.5.5}\label{eq:2.5.5}
\end{equation}
where $N(p)$ is a numerical normalization factor, to be chosen later. Up to
this point, we have said nothing about how the $\sigma$ labels are related for
different momenta; Eq.~\eqref{eq:2.5.5} now fills that gap.

Operating on \eqref{eq:2.5.5} with an arbitrary homogeneous Lorentz transformation
$U(\Lambda)$, we now find
\begin{equation}
\begin{aligned}
U(\Lambda)\ket{p,\sigma}
&=N(p)U(\Lambda L(p))\ket{k,\sigma}\\
&=N(p)U(L(\Lambda p))
U(L^{-1}(\Lambda p)\Lambda L(p))\ket{k,\sigma}\,.
\end{aligned}
\tag{2.5.6}\label{eq:2.5.6}
\end{equation}
The point of this last step is that the Lorentz transformation
$L^{-1}(\Lambda p)\Lambda L(p)$ takes $k$ to $L(p)k=p$, and then to
$\Lambda p$, and then back to $k$, so it belongs to the subgroup of the
homogeneous Lorentz group consisting of Lorentz transformations
$W^\mu{}_{\nu}$ that leave $k^\mu$ invariant:
\begin{equation}
W^\mu{}_{\nu}k^\nu=k^\mu\,.
\tag{2.5.7}\label{eq:2.5.7}
\end{equation}
This subgroup is called the \emph{little group}~\hyperref[ref:2.5]{[5]}. For any
$W$ satisfying Eq.~\eqref{eq:2.5.7}, we have
\begin{equation}
U(W)\ket{k,\sigma}
=\sum_{\sigma'}D_{\sigma'\sigma}(W)\ket{k,\sigma'}\,.
\tag{2.5.8}\label{eq:2.5.8}
\end{equation}
The coefficients $D(W)$ furnish a representation of the little group; that is,
for any elements $\bar W,W$ we have
\[
\begin{aligned}
\sum_{\sigma'}D_{\sigma'\sigma}(\bar W W)\ket{k,\sigma'}
&=U(\bar W W)\ket{k,\sigma}=U(\bar W)U(W)\ket{k,\sigma}\\
&=U(\bar W)\sum_{\sigma''}D_{\sigma''\sigma}(W)\ket{k,\sigma''}\\
&=\sum_{\sigma'\sigma''}D_{\sigma'\sigma''}(\bar W)
D_{\sigma''\sigma}(W)\ket{k,\sigma'}
\end{aligned}
\]
and so
\begin{equation}
D_{\sigma'\sigma}(\bar W W)
=\sum_{\sigma''}D_{\sigma'\sigma''}(\bar W)D_{\sigma''\sigma}(W)\,.
\tag{2.5.9}\label{eq:2.5.9}
\end{equation}

In particular, we may apply Eq.~\eqref{eq:2.5.8} to the little-group transformation
\begin{equation}
W(\Lambda,p)\equiv L^{-1}(\Lambda p)\Lambda L(p)
\tag{2.5.10}\label{eq:2.5.10}
\end{equation}
and then Eq.~\eqref{eq:2.5.6} takes the form
\[
U(\Lambda)\ket{p,\sigma}
=N(p)\sum_{\sigma'}D_{\sigma'\sigma}(W(\Lambda,p))
U(L(\Lambda p))\ket{k,\sigma'}\,,
\]
or, recalling the definition \eqref{eq:2.5.5}:
\begin{equation}
U(\Lambda)\ket{p,\sigma}
=\left(\frac{N(p)}{N(\Lambda p)}\right)
\sum_{\sigma'}D_{\sigma'\sigma}(W(\Lambda,p))
\ket{\Lambda p,\sigma'}\,.
\tag{2.5.11}\label{eq:2.5.11}
\end{equation}
Apart from the question of normalization, the problem of determining the
coefficients $C_{\sigma'\sigma}$ in the transformation rule \eqref{eq:2.5.3} has been
reduced to the problem of finding the representations of the little group.
This approach, of deriving representations of a group like the
Poincar\'e group from the representations of a little group, is called the method
of induced representations~\hyperref[ref:2.6]{[6]}.

Table 2.1 gives a convenient choice of the standard momentum $k^\mu$ and the
corresponding little group for the various classes of four-momenta.

Of these six classes of four-momenta, only (a), (c), and (f) have any known
interpretations in terms of physical states. Not much needs to be said here
about case (f)---$p^\mu=0$; it describes the vacuum, which is simply left
invariant by $U(\Lambda)$. In what follows we will consider only cases (a) and
(c), which cover particles of mass $M>0$ and mass zero, respectively.

This is a good place to pause, and say something about the normalization of
these states. By the usual orthonormalization procedure of quantum mechanics,
we may choose the states with standard momentum $k^\mu$ to be orthonormal, in
the sense that
\begin{equation}
\braket{k',\sigma'}{k,\sigma}
=\delta^3(\mathbf{k}'-\mathbf{k})\delta_{\sigma'\sigma}\,.
\tag{2.5.12}\label{eq:2.5.12}
\end{equation}
(The delta function appears here because $\ket{k,\sigma}$ and
$\ket{k',\sigma'}$ are eigenstates of a Hermitian operator with eigenvalues
$\mathbf{k}$ and $\mathbf{k}'$, respectively.) This has the immediate
consequence that the representation of the little group in Eqs.~\eqref{eq:2.5.8} and
\eqref{eq:2.5.11} must be unitary\footnote{The little groups $SO(2,1)$ and $SO(3,1)$
for $p^2>0$ and $p^\mu=0$ have no non-trivial finite-dimensional unitary
representations, so if there were any states with a given momentum $p^\mu$
with $p^2>0$ or $p^\mu=0$ that transform non-trivially under the little group,
there would have to be an infinite number of them.}
\begin{equation}
D^\dagger(W)=D^{-1}(W)\,.
\tag{2.5.13}\label{eq:2.5.13}
\end{equation}

\begin{center}
\begin{minipage}{0.94\textwidth}
\noindent Table 2.1.\quad Standard momenta and the corresponding little group for various
classes of four-momenta. Here $\kappa$ is an arbitrary positive energy, say
1 eV. The little groups are mostly pretty obvious: $SO(3)$ is the ordinary
rotation group in three dimensions (excluding space inversions), because
rotations are the only proper orthochronous Lorentz transformations that leave
at rest a particle with zero momentum, while $SO(2,1)$ and $SO(3,1)$ are the
Lorentz groups in $(2+1)$- and $(3+1)$-dimensions, respectively. The group
$ISO(2)$ is the group of Euclidean geometry, consisting of rotations and
translations in two dimensions. Its appearance as the little group for
$p^2=0$ is explained below.

\medskip
\centering
\renewcommand{\arraystretch}{1.35}
\begin{tabular}{@{}lll@{}}
\hline\hline
 & Standard $k^\mu$ & Little Group \\
\hline
(a) $p^2=-M^2<0,\ p^0>0$ & $(M,0,0,0)$ & $SO(3)$ \\
(b) $p^2=-M^2<0,\ p^0<0$ & $(-M,0,0,0)$ & $SO(3)$ \\
(c) $p^2=0,\ p^0>0$ & $(\kappa,0,0,\kappa)$ & $ISO(2)$ \\
(d) $p^2=0,\ p^0<0$ & $(-\kappa,0,0,\kappa)$ & $ISO(2)$ \\
(e) $p^2=N^2>0$ & $(0,0,0,N)$ & $SO(2,1)$ \\
(f) $p^\mu=0$ & $(0,0,0,0)$ & $SO(3,1)$ \\
\hline\hline
\end{tabular}
\end{minipage}
\end{center}

Now, what about scalar products for arbitrary momenta? Using the unitarity of
the operator $U(\Lambda)$ in Eqs.~\eqref{eq:2.5.5} and \eqref{eq:2.5.11}, we find for the scalar
product:
\[
\begin{aligned}
\braket{p',\sigma'}{p,\sigma}
&=N(p)\braket{U^{-1}(L(p))p',\sigma'}{k,\sigma}\\
&=N(p)N^*(p')
D_{\sigma\sigma'}(W(L^{-1}(p),p'))^*
\delta^3(\mathbf{k}'-\mathbf{k})\,,
\end{aligned}
\]
where $k'\equiv L^{-1}(p)p'$. Since also $k=L^{-1}(p)p$, the delta function
$\delta^3(\mathbf{k}-\mathbf{k}')$ is proportional to
$\delta^3(\mathbf{p}-\mathbf{p}')$. For $p'=p$, the little-group
transformation here is trivial, $W(L^{-1}(p),p)=\mathbf{1}$, and so the scalar product is
\begin{equation}
\braket{p',\sigma'}{p,\sigma}
=|N(p)|^2\delta_{\sigma'\sigma}
\delta^3(\mathbf{k}'-\mathbf{k})\,.
\tag{2.5.14}\label{eq:2.5.14}
\end{equation}

It remains to work out the proportionality factor relating
$\delta^3(\mathbf{k}-\mathbf{k}')$ and
$\delta^3(\mathbf{p}-\mathbf{p}')$. Note that the Lorentz-invariant integral
of an arbitrary scalar function $f(p)$ over four-momenta with
$-p^2=M^2\geq0$ and $p^0>0$ (i.e., cases (a) or (c)) may be written
\[
\begin{aligned}
\int d^4p\,\delta(p^2+M^2)\theta(p^0)f(p)
&=\int d^3\mathbf{p}\,dp^0\,
\delta((p^0)^2-\mathbf{p}^2-M^2)\theta(p^0)f(\mathbf{p},p^0)\\
&=\int d^3\mathbf{p}\,
\frac{f(\mathbf{p},\sqrt{\mathbf{p}^2+M^2})}
{2\sqrt{\mathbf{p}^2+M^2}}\,.
\end{aligned}
\]
($\theta(p^0)$ is the step function: $\theta(x)=1$ for $x\geq0$,
$\theta(x)=0$ for $x<0$.) We see that when integrating on the `mass shell'
$p^2+M^2=0$, the invariant volume element is
\begin{equation}
d^3\mathbf{p}/\sqrt{\mathbf{p}^2+M^2}\,.
\tag{2.5.15}\label{eq:2.5.15}
\end{equation}
The delta function is defined by
\[
\begin{aligned}
F(\mathbf{p})
&=\int F(\mathbf{p}')\delta^3(\mathbf{p}-\mathbf{p}')d^3\mathbf{p}'\\
&=\int F(\mathbf{p}')
\left[\sqrt{\mathbf{p}'^2+M^2}\,
\delta^3(\mathbf{p}'-\mathbf{p})\right]
\frac{d^3\mathbf{p}'}{\sqrt{\mathbf{p}'^2+M^2}}\,,
\end{aligned}
\]
so we see that the invariant delta function is
\begin{equation}
\sqrt{\mathbf{p}'^2+M^2}\,\delta^3(\mathbf{p}'-\mathbf{p})
=p^0\delta^3(\mathbf{p}'-\mathbf{p})\,.
\tag{2.5.16}\label{eq:2.5.16}
\end{equation}
Since $p'$ and $p$ are related to $k'$ and $k$ respectively by a Lorentz
transformation, $L(p)$, we have then
\[
p^0\delta^3(\mathbf{p}'-\mathbf{p})
=k^0\delta^3(\mathbf{k}'-\mathbf{k})
\]
and therefore
\begin{equation}
\braket{p',\sigma'}{p,\sigma}
=|N(p)|^2\delta_{\sigma'\sigma}
\left(\frac{p^0}{k^0}\right)
\delta^3(\mathbf{p}'-\mathbf{p})\,.
\tag{2.5.17}\label{eq:2.5.17}
\end{equation}
The normalization factor $N(p)$ is sometimes chosen to be just $N(p)=1$, but
then we would need to keep track of the $p^0/k^0$ factor in scalar products.
Instead, I will here adopt the more usual convention that
\begin{equation}
N(p)=\sqrt{k^0/p^0}
\tag{2.5.18}\label{eq:2.5.18}
\end{equation}
for which
\begin{equation}
\braket{p',\sigma'}{p,\sigma}
=\delta_{\sigma'\sigma}\delta^3(\mathbf{p}'-\mathbf{p})\,.
\tag{2.5.19}\label{eq:2.5.19}
\end{equation}
We now consider the two cases of physical interest: particles of mass $M>0$,
and particles of zero mass.

\subsubsection*{Mass Positive-Definite}

The little group here is the three-dimensional rotation group. Its unitary
representations can be broken up into a direct sum of irreducible unitary
representations~\hyperref[ref:2.7]{[7]} $D^{(j)}_{\sigma'\sigma}(R)$ of dimensionality
$2j+1$, with $j=0,\tfrac12,1,\ldots$. These can be built up from the standard
matrices for infinitesimal rotations $R_{ik}=\delta_{ik}+\Theta_{ik}$, with
$\Theta_{ik}=-\Theta_{ki}$ infinitesimal:
\begin{equation}
D^{(j)}_{\sigma'\sigma}(\mathbf{1}+\Theta)
=\delta_{\sigma'\sigma}
+\frac{i}{2}\Theta_{ik}(J^{(j)}_{ik})_{\sigma'\sigma}\,,
\tag{2.5.20}\label{eq:2.5.20}
\end{equation}
\begin{equation}
\begin{aligned}
(J^{(j)}_{23}\pm iJ^{(j)}_{31})_{\sigma'\sigma}
&=(J^{(j)}_1\pm iJ^{(j)}_2)_{\sigma'\sigma}\\
&=\delta_{\sigma',\sigma\pm1}
\sqrt{(j\mp\sigma)(j\pm\sigma+1)}\,,
\end{aligned}
\tag{2.5.21}\label{eq:2.5.21}
\end{equation}
\begin{equation}
(J^{(j)}_{12})_{\sigma'\sigma}
=(J^{(j)}_3)_{\sigma'\sigma}
=\sigma\delta_{\sigma'\sigma}\,,
\tag{2.5.22}\label{eq:2.5.22}
\end{equation}
with $\sigma$ running over the values $j,j-1,\ldots,-j$. For a particle of
mass $M>0$ and spin $j$, Eq.~\eqref{eq:2.5.11} now becomes
\begin{equation}
U(\Lambda)\ket{p,\sigma}
=\sqrt{\frac{(\Lambda p)^0}{p^0}}
\sum_{\sigma'}D^{(j)}_{\sigma'\sigma}(W(\Lambda,p))
\ket{\Lambda p,\sigma'}\,,
\tag{2.5.23}\label{eq:2.5.23}
\end{equation}
with the little-group element $W(\Lambda,p)$ (the Wigner
rotation~\hyperref[ref:2.5]{[5]}) given by Eq.~\eqref{eq:2.5.10}:
\[
W(\Lambda,p)=L^{-1}(\Lambda p)\Lambda L(p)\,.
\]
To calculate this rotation, we need to choose a `standard boost' $L(p)$ which
carries the four-momentum from $k^\mu=(M,\mathbf{0})$ to $p^\mu$. This is
conveniently chosen as
\begin{equation}
\begin{aligned}
L^i{}_k(p)&=\delta_{ik}+(\gamma-1)\hat p_i\hat p_k\,,\\
L^i{}_0(p)&=L^0{}_i(p)=\hat p_i\sqrt{\gamma^2-1}\,,\\
L^0{}_0(p)&=\gamma\,,
\end{aligned}
\tag{2.5.24}\label{eq:2.5.24}
\end{equation}
where
\[
\hat p_i\equiv p_i/|\mathbf{p}|\,,
\qquad
\gamma\equiv\sqrt{\mathbf{p}^2+M^2}/M\,.
\]
It is very important that when $\Lambda^\mu{}_{\nu}$ is an arbitrary
three-dimensional rotation $\mathcal{R}$, the Wigner rotation
$W(\Lambda,p)$ is the same as $\mathcal{R}$ for all $p$. To see this, note that
the boost \eqref{eq:2.5.24} may be expressed as
\[
L(p)=R(\hat{\mathbf{p}})B(|\mathbf{p}|)R^{-1}(\hat{\mathbf{p}})\,,
\]
where $R(\hat{\mathbf{p}})$ is a rotation (to be defined in a standard way
below, in Eq.~\eqref{eq:2.5.47}) that takes the three-axis into the direction of
$\mathbf{p}$, and
\[
B(|\mathbf{p}|)=
\begin{pmatrix}
\gamma&0&0&\sqrt{\gamma^2-1}\\
0&1&0&0\\
0&0&1&0\\
\sqrt{\gamma^2-1}&0&0&\gamma
\end{pmatrix}.
\]
Then for an arbitrary rotation $\mathcal{R}$
\[
W(\mathcal{R},p)=R(\mathcal{R}\hat{\mathbf{p}})
B^{-1}(|\mathbf{p}|)R^{-1}(\mathcal{R}\hat{\mathbf{p}})
\mathcal{R}R(\hat{\mathbf{p}})B(|\mathbf{p}|)R^{-1}(\hat{\mathbf{p}})\,.
\]
But the rotation
$R^{-1}(\mathcal{R}\hat{\mathbf{p}})\mathcal{R}R(\hat{\mathbf{p}})$ takes the
three-axis into the direction $\hat{\mathbf{p}}$, and then into the direction
$\mathcal{R}\hat{\mathbf{p}}$, and then back to the three-axis, so it must be
just a rotation by some angle $\theta$ around the three-axis
\[
R^{-1}(\mathcal{R}\hat{\mathbf{p}})\mathcal{R}R(\hat{\mathbf{p}})
=R(\theta)\equiv
\begin{pmatrix}
1&0&0&0\\
0&\cos\theta&\sin\theta&0\\
0&-\sin\theta&\cos\theta&0\\
0&0&0&1
\end{pmatrix}.
\]
Since $R(\theta)$ commutes with $B(|\mathbf{p}|)$, this now gives
\[
\begin{aligned}
W(\mathcal{R},p)
&=R(\mathcal{R}\hat{\mathbf{p}})B^{-1}(|\mathbf{p}|)R(\theta)
B(|\mathbf{p}|)R^{-1}(\hat{\mathbf{p}})\\
&=R(\mathcal{R}\hat{\mathbf{p}})R(\theta)R^{-1}(\hat{\mathbf{p}})
\end{aligned}
\]
and hence
\[
W(\mathcal{R},p)=\mathcal{R}
\]
as was to be shown. Thus states of a moving massive particle (and, by
extension, multi-particle states) have the same transformation under rotations
as in non-relativistic quantum mechanics. This is another piece of good
news---the whole apparatus of spherical harmonics, Clebsch--Gordan
coefficients, etc. can be carried over wholesale from non-relativistic to
relativistic quantum mechanics.

\subsubsection*{Mass Zero}

First, we have to work out the structure of the little group. Consider an
arbitrary little-group element $W^\mu{}_{\nu}$, with
$W^\mu{}_{\nu}k^\nu=k^\mu$, where $k^\mu$ is the standard four-momentum for
this case, $k^\mu=(1,0,0,1)$. Acting on a time-like four-vector
$t^\mu=(1,0,0,0)$, such a Lorentz transformation must yield a four-vector
$Wt$ whose length and scalar product with $Wk=k$ are the same as those of $t$:
\[
(Wt)^\mu(Wt)_\mu=t^\mu t_\mu=-1\,,
\qquad
(Wt)^\mu k_\mu=t^\mu k_\mu=-1\,.
\]
Any four-vector that satisfies the second condition may be written
\[
(Wt)^\mu=(1+\zeta,\alpha,\beta,\zeta)
\]
and the first condition then yields the relation
\begin{equation}
\zeta=(\alpha^2+\beta^2)/2\,.
\tag{2.5.25}\label{eq:2.5.25}
\end{equation}
It follows that the effect of $W^\mu{}_{\nu}$ on $t^\nu$ is the same as that of
the Lorentz transformation
\begin{equation}
S^\mu{}_{\nu}(\alpha,\beta)=
\begin{pmatrix}
1+\zeta&\alpha&\beta&-\zeta\\
\alpha&1&0&-\alpha\\
\beta&0&1&-\beta\\
\zeta&\alpha&\beta&1-\zeta
\end{pmatrix}.
\tag{2.5.26}\label{eq:2.5.26}
\end{equation}
This does not mean that $W$ equals $S(\alpha,\beta)$, but it does mean that
$S^{-1}(\alpha,\beta)W$ is a Lorentz transformation that leaves the time-like
four-vector $(1,0,0,0)$ invariant, and is therefore a pure rotation. Also,
$S^\mu{}_{\nu}$ like $W^\mu{}_{\nu}$ leaves the light-like four-vector
$(1,0,0,1)$ invariant, so $S^{-1}(\alpha,\beta)W$ must be a rotation by some
angle $\theta$ around the three-axis
\begin{equation}
S^{-1}(\alpha,\beta)W=R(\theta)\,,
\tag{2.5.27}\label{eq:2.5.27}
\end{equation}
where
\[
R^\mu{}_{\nu}(\theta)\equiv
\begin{pmatrix}
1&0&0&0\\
0&\cos\theta&\sin\theta&0\\
0&-\sin\theta&\cos\theta&0\\
0&0&0&1
\end{pmatrix}.
\]
The most general element of the little group is therefore of the form
\begin{equation}
W(\theta,\alpha,\beta)=S(\alpha,\beta)R(\theta)\,.
\tag{2.5.28}\label{eq:2.5.28}
\end{equation}
What group is this? We note that the transformations with $\theta=0$ or with
$\alpha=\beta=0$ form subgroups:
\begin{equation}
S(\bar\alpha,\bar\beta)S(\alpha,\beta)
=S(\bar\alpha+\alpha,\bar\beta+\beta)
\tag{2.5.29}\label{eq:2.5.29}
\end{equation}
\begin{equation}
R(\bar\theta)R(\theta)=R(\bar\theta+\theta)\,.
\tag{2.5.30}\label{eq:2.5.30}
\end{equation}
These subgroups are \emph{Abelian}---that is, their elements all commute with
each other. Furthermore, the subgroup with $\theta=0$ is \emph{invariant}, in
the sense that its elements are transformed into other elements of the same
subgroup by any member of the group
\begin{equation}
R(\theta)S(\alpha,\beta)R^{-1}(\theta)
=S(\alpha\cos\theta+\beta\sin\theta,
-\alpha\sin\theta+\beta\cos\theta)\,.
\tag{2.5.31}\label{eq:2.5.31}
\end{equation}
From Eqs.~\eqref{eq:2.5.29}--\eqref{eq:2.5.31} we can work out the product of any group elements.
The reader will recognize these multiplication rules as those of the group
$ISO(2)$, consisting of translations (by a vector $(\alpha,\beta)$) and
rotations (by an angle $\theta$) in two dimensions.

Groups that do \emph{not} have invariant Abelian subgroups have certain simple
properties, and for this reason are called \emph{semi-simple}. As we have seen,
the little group $ISO(2)$ like the Poincar\'e group is \emph{not}
semi-simple, and this leads to interesting complications. First, let's take a
look at the Lie algebra of $ISO(2)$. For $\theta,\alpha,\beta$ infinitesimal,
the general group element is
\[
W(\theta,\alpha,\beta)^\mu{}_{\nu}
=\delta^\mu{}_{\nu}+\omega^\mu{}_{\nu}\,,
\qquad
\omega_{\mu\nu}=
\begin{pmatrix}
0&-\alpha&-\beta&0\\
\alpha&0&\theta&-\alpha\\
\beta&-\theta&0&-\beta\\
0&\alpha&\beta&0
\end{pmatrix}.
\]
From \eqref{eq:2.4.3}, we see then that the corresponding Hilbert space operator is
\begin{equation}
U(W(\theta,\alpha,\beta))=\mathbf{1}+i\alpha A+i\beta B+i\theta J_3\,,
\tag{2.5.32}\label{eq:2.5.32}
\end{equation}
where $A$ and $B$ are the Hermitian operators
\begin{equation}
A=-J^{13}+J^{10}=J_2+K_1\,,
\tag{2.5.33}\label{eq:2.5.33}
\end{equation}
\begin{equation}
B=-J^{23}+J^{20}=-J_1+K_2\,,
\tag{2.5.34}\label{eq:2.5.34}
\end{equation}
and, as before, $J_3=J_{12}$. Either from \eqref{eq:2.4.18}--\eqref{eq:2.4.20}, or directly from
Eqs.~\eqref{eq:2.5.29}--\eqref{eq:2.5.31}, we see that these generators have the commutators
\begin{align}
[J_3,A]&=+iB\,, \tag{2.5.35}\label{eq:2.5.35}\\
[J_3,B]&=-iA\,, \tag{2.5.36}\label{eq:2.5.36}\\
[A,B]&=0\,. \tag{2.5.37}\label{eq:2.5.37}
\end{align}
Since $A$ and $B$ are commuting Hermitian operators they (like the momentum
generators of the Poincar\'e group) can be simultaneously
diagonalized by states $\ket{k,a,b}$
\[
A\ket{k,a,b}=a\ket{k,a,b}\,,
\qquad
B\ket{k,a,b}=b\ket{k,a,b}\,.
\]
The problem is that if we find one such set of non-zero eigenvalues of $A,B$,
then we find a whole continuum. From Eq.~\eqref{eq:2.5.31}, we have
\[
U[R(\theta)]A\,U^{-1}[R(\theta)]=A\cos\theta-B\sin\theta\,,
\]
\[
U[R(\theta)]B\,U^{-1}[R(\theta)]=A\sin\theta+B\cos\theta\,,
\]
and so, for arbitrary $\theta$,
\[
A\ket{k,a,b}^{\theta}=(a\cos\theta-b\sin\theta)\ket{k,a,b}^{\theta}\,,
\]
\[
B\ket{k,a,b}^{\theta}=(a\sin\theta+b\cos\theta)\ket{k,a,b}^{\theta}\,,
\]
where
\[
\ket{k,a,b}^{\theta}\equiv U^{-1}(R(\theta))\ket{k,a,b}\,.
\]

Massless particles are not observed to have any continuous degree of freedom
like $\theta$; to avoid such a continuum of states, we must require that
physical states (now called $\ket{k,\sigma}$) are eigenvectors of $A$ and $B$
with $a=b=0$:
\begin{equation}
A\ket{k,\sigma}=B\ket{k,\sigma}=0\,.
\tag{2.5.38}\label{eq:2.5.38}
\end{equation}
These states are then distinguished by the eigenvalue of the remaining
generator
\begin{equation}
J_3\ket{k,\sigma}=\sigma\ket{k,\sigma}\,.
\tag{2.5.39}\label{eq:2.5.39}
\end{equation}
Since the momentum $\mathbf{k}$ is in the three-direction, $\sigma$ gives the
component of angular momentum in the direction of motion, or \emph{helicity}.

We are now in a position to calculate the Lorentz transformation properties of
general massless particle states. First note that by use of the general
arguments of Section 2.2, Eq.~\eqref{eq:2.5.32} generalizes for finite $\alpha$ and
$\beta$ to
\begin{equation}
U(S(\alpha,\beta))=\exp(i\alpha A+i\beta B)
\tag{2.5.40}\label{eq:2.5.40}
\end{equation}
and for finite $\theta$ to
\begin{equation}
U(R(\theta))=\exp(iJ_3\theta)\,.
\tag{2.5.41}\label{eq:2.5.41}
\end{equation}
An arbitrary element $W$ of the little group can be put in the form \eqref{eq:2.5.28},
so that
\[
U(W)\ket{k,\sigma}
=\exp(i\alpha A+i\beta B)\exp(i\theta J_3)\ket{k,\sigma}
=\exp(i\theta\sigma)\ket{k,\sigma}
\]
and therefore Eq.~\eqref{eq:2.5.8} gives
\[
D_{\sigma'\sigma}(W)=\exp(i\theta\sigma)\delta_{\sigma'\sigma}\,,
\]
where $\theta$ is the angle defined by expressing $W$ as in Eq.~\eqref{eq:2.5.28}. The
Lorentz transformation rule for a massless particle of arbitrary helicity is
now given by Eqs.~\eqref{eq:2.5.11} and \eqref{eq:2.5.18} as
\begin{equation}
U(\Lambda)\ket{p,\sigma}
=\sqrt{\frac{(\Lambda p)^0}{p^0}}
\exp(i\sigma\theta(\Lambda,p))\ket{\Lambda p,\sigma}
\tag{2.5.42}\label{eq:2.5.42}
\end{equation}
with $\theta(\Lambda,p)$ defined by
\begin{equation}
W(\Lambda,p)\equiv L^{-1}(\Lambda p)\Lambda L(p)
\equiv S(\alpha(\Lambda,p),\beta(\Lambda,p))R(\theta(\Lambda,p))\,.
\tag{2.5.43}\label{eq:2.5.43}
\end{equation}
We shall see in Section 5.9 that electromagnetic gauge invariance arises from
the part of the little group parameterized by $\alpha$ and $\beta$.

At this point we have not yet encountered any reason that would forbid the
helicity $\sigma$ of a massless particle from being an arbitrary real number.
As we shall see in Section 2.7, there are topological considerations that
restrict the allowed values of $\sigma$ to integers and half-integers, just as
for massive particles.

To calculate the little-group element \eqref{eq:2.5.43} for a given $\Lambda$ and $p$,
(and also to enable us to calculate the effect of space or time inversion on
these states in the next section) we need to fix a convention for the standard
Lorentz transformation that takes us from
$k^\mu=(\kappa,0,0,\kappa)$ to $p^\mu$. This may conveniently be chosen to have
the form
\begin{equation}
L(p)=R(\hat{\mathbf{p}})B(|\mathbf{p}|/\kappa)
\tag{2.5.44}\label{eq:2.5.44}
\end{equation}
where $B(u)$ is a pure boost along the three-direction:
\begin{equation}
B(u)\equiv
\begin{pmatrix}
(u^2+1)/2u&0&0&(u^2-1)/2u\\
0&1&0&0\\
0&0&1&0\\
(u^2-1)/2u&0&0&(u^2+1)/2u
\end{pmatrix}
\tag{2.5.45}\label{eq:2.5.45}
\end{equation}
and $R(\hat{\mathbf{p}})$ is a pure rotation that carries the three-axis into
the direction of the unit vector $\hat{\mathbf{p}}$. For instance, suppose we
take $\hat{\mathbf{p}}$ to have polar and azimuthal angles $\theta$ and $\phi$:
\begin{equation}
\hat{\mathbf{p}}=(\sin\theta\cos\phi,\sin\theta\sin\phi,\cos\theta)\,.
\tag{2.5.46}\label{eq:2.5.46}
\end{equation}
Then we can take $R(\hat{\mathbf{p}})$ as a rotation by angle $\theta$ around
the two-axis, which takes $(0,0,1)$ into $(\sin\theta,0,\cos\theta)$, followed
by a rotation by angle $\phi$ around the three-axis:
\begin{equation}
U(R(\hat{\mathbf{p}}))=\exp(i\phi J_3)\exp(i\theta J_2)\,,
\tag{2.5.47}\label{eq:2.5.47}
\end{equation}
where $0\leq\theta\leq\pi$, $0\leq\phi<2\pi$. (We give
$U(R(\hat{\mathbf{p}}))$ rather than $R(\hat{\mathbf{p}})$, together with a
specification of the range of $\phi$ and $\theta$, because shifting $\theta$
or $\phi$ by $2\pi$ would give the same rotation $R(\hat{\mathbf{p}})$, but a
different sign for $U(R(\hat{\mathbf{p}}))$ when acting on half-integer spin
states.) Since \eqref{eq:2.5.47} is a rotation, and does take the three-axis into the
direction \eqref{eq:2.5.46}, any other choice of such an $R(\hat{\mathbf{p}})$ would
differ from this one by at most an initial rotation around the three-axis,
corresponding to a mere redefinition of the phase of the one-particle states.

Note that the helicity is Lorentz-invariant; a massless particle of a given
helicity $\sigma$ looks the same (aside from its momentum) in all inertial
frames. Indeed, we would be justified in thinking of massless particles of
each different helicity as different species of particles. However, as we
shall see in the next section, particles of opposite helicity are related by
the symmetry of space inversion. Thus, because electromagnetic and
gravitational forces obey space inversion symmetry, the massless particles of
helicity $\pm1$ associated with electromagnetic phenomena are both called
\emph{photons}, and the massless particles of helicity $\pm2$ that are believed
to be associated with gravitation are both called \emph{gravitons}. On the
other hand, the supposedly massless particles of helicity $\pm1/2$ that are
emitted in nuclear beta decay have no interactions (apart from gravitation)
that respect the symmetry of space inversion, so these particles are given
different names: \emph{neutrinos} for helicity $+1/2$, and
\emph{antineutrinos} for helicity $-1/2$.

Even though the helicity of a massless particle is Lorentz-invariant, the state
itself is not. In particular, because of the helicity-dependent phase factor
$\exp(i\sigma\theta)$ in Eq.~\eqref{eq:2.5.42}, a state formed as a linear superposition
of one-particle states with opposite helicities will be changed by a Lorentz
transformation into a different superposition. For instance, a general
one-photon state of four-momenta may be written
\[
\ket{p;\alpha}=\alpha_+\ket{p,+1}+\alpha_-\ket{p,-1}\,,
\]
where
\[
|\alpha_+|^2+|\alpha_-|^2=1\,.
\]
The generic case is one of \emph{elliptic polarization}, with
$|\alpha_\pm|$ both non-zero and unequal. \emph{Circular polarization} is the
limiting case where either $\alpha_+$ or $\alpha_-$ vanishes, and
\emph{linear polarization} is the opposite extreme, with
$|\alpha_+|=|\alpha_-|$. The overall phase of $\alpha_+$ and $\alpha_-$ has no
physical significance, and for linear polarization may be adjusted so that
$\alpha_-=\alpha_+^*$, but the relative phase is still important. Indeed, for
linear polarizations with $\alpha_-=\alpha_+^*$, the phase of $\alpha_+$ may be
identified as the angle between the plane of polarization and some fixed
reference direction perpendicular to $\mathbf{p}$. Eq.~\eqref{eq:2.5.42} shows that
under a Lorentz transformation $\Lambda^\mu{}_{\nu}$, this angle rotates by an
amount $\theta(\Lambda,p)$. Plane polarized gravitons can be defined in a
similar way, and here Eq.~\eqref{eq:2.5.42} has the consequence that a Lorentz
transformation $\Lambda$ rotates the plane of polarization by an angle
$2\theta(\Lambda,p)$.
